\begin{statementbox}
	\textbf{\textcolor{CStatement}{条件}}
	\begin{description}[style=nextline, leftmargin=2.8em, labelsep=0.5em, itemsep=0.35em]
		\item[\textbf{\textcolor{CStatement}{条件1（直径条件）：}}] $AB$ 是 $\odot O$ 的直径。
		\item[\textbf{\textcolor{CStatement}{条件2（点与端点）：}}] 点 $C$ 在 $\odot O$ 上且 $C\neq A,B$。
	\end{description}

	\textbf{\textcolor{CStatement}{结论}}
	\begin{description}[style=nextline, leftmargin=2.8em, labelsep=0.5em, itemsep=0.45em]
		\item[\textbf{\textcolor{CStatement}{结论一（圆周角直角）：}}]
			\textbf{\textcolor{CStatement}{直观描述：}}
			直径所对的圆周角等于 $90^\circ$。
			\[
				\angle ACB = 90^\circ.
			\]

		\item[\textbf{\textcolor{CStatement}{推广结论（逆定理）：}}]
			\textbf{\textcolor{CStatement}{直观描述：}}
			若圆周角为 $90^\circ$，则其所对的弦过圆心（即为直径）。
			\[
				\angle ACB = 90^\circ \implies O\in AB.
			\]
	\end{description}
\end{statementbox}
